% r03_h1s_e: proof of the two-flow comparison bound, part 5: the telescoping, every term is small, the exponent theta_1, the sectors, the proof of the theorem.

\paragraph{The telescoping.}
After (NF1)--(NF4) and (R1)--(R3) in each flow's own frame, $E^F(\cR)$ for one ordering class is a finite sum of multilinear expressions in the elements: slice lines and top pieces (with at most two lattice derivatives, \Cref{rot:h1s:nd}), compounds $(U,f)$ and $(U_1,l,U_2)$ (own-frame remainder sums, \Cref{rot:h1s:U}), first-order compounds $(T,f)$, $(SC,f)$, the brackets $[G(t)-G(x)]$ of (R2b) with one differenced factor, the coarse elements $\Gamma^F$, $H^F$, $B^F$, the externals and vertex functions at $\tau$-matched points (shifts included), the contour elements at $e>0$, and the common numbers (the zero moments $c$, the order-four shares, the vacuum-polarisation leftover weight, the one-point structures), which are never switched. The two flows' expressions have the same combinatorics. Order the elements consistently (compounds first, then lines, coarse elements, externals, contour elements) and telescope:
\[
E^X(\cR)-E^Y(\cR)=\sum_iT_i,\qquad T_i\defeq\big[e_1,\dots,e_{i-1}\text{ in the $Y\circ\tau$ version};\ e_i\to e_i^X-e_i^Y\circ\tau;\ e_{i+1},\dots\text{ in the $X$ version}\big],
\]
an identity, every term being the same multilinear form with elements substituted, the co-switched pairs (a self-energy and its counterterm share, the two Wick contractions of \Cref{rot:h1s:M}, a scalar W and its order-four share) switched together; at most $n_\cR$ terms per class, and finitely many classes per family. Compounds are switched first so that a unit and its leg are compared as a whole and no mixed-frame first moment is ever written; externals are switched late so that the only frame-sensitive quantity left is the difference of the two flows' own-frame remainders, which \Cref{rot:h1s:E} bounds. Identities (the counterterm cancellations, the Ward sum, the moment subtractions) are used only at the two pure endpoints, in each flow's own frame; the mixed terms are bounded in absolute value.

\begin{lemma}[every telescoping term is small]\label{rot:h1s:T}
For every ordering class and every term $T_i$, with $\theta\le\theta_1=\alpha^2/(4(1+\alpha)^2)$,
\[
|T_i|\le C(1+k)^p\eta'^\theta\cdot[\text{the per-flow bound \cite[(1.33)]{Ba3} of the monomial}],\qquad p=p(\bar n),
\]
with the prefactor $p_r^2\delta_r^{-1/2}$ when $e_i$ is a contour element.
\end{lemma}
\begin{proof}
Apply \Cref{rot:h1s:MS} to $|T_i|$, the absolute value of a multilinear form whose degrees are those of \Cref{rot:h1s:Dp}(a).
(i) \emph{Per-flow bounds} of every element in $X$-coordinates, whichever flow it comes from: lines by \eqref{eq:h1s:B}, with $\delta/2$ for the $Y\circ\tau$ versions; compounds by \Cref{rot:h1s:U}(i),(iii), lines of dimension $-2-a$, resp.\ $-3-a$ ($0$ at the top scale for a $\Gamma$ middle), at their leg's scale, or vertex weights $\alpha$ for external legs; first-order compounds by \Cref{rot:h1s:U}(iv) (degree-one objects); the brackets of (R2b) by $|G(t)-G(x)|\le|t-x|\sup|\nabla g|$ plus the explicit-$\eta'$ displacement term; coarse elements, externals and vertex functions by their table; the common numbers by \Cref{rot:h1s:Z0} and \Cref{rot:h1s:MS} per flow ($\le C(1+r)^{n_0+1}$ each; the vacuum-polarisation leftover $\le C\ell_J^{-1}$ as a weight $-1$); at $e>0$ the contour elements by \Cref{rot:prop:BD}(b), a contour line that received a derivative being a line of dimension $-1$ at its slice (the contour sum carries the weight $\eta$ per bond against $\sim\ell_j/\eta$ bonds within the decay range, \cite[(3.63)--(3.64)]{K1} with the contour clause of \cite[Prop.~3.7]{K1}). Decay is re-anchored at merged vertices as in the first step of \Cref{rot:h1s:MS}.
(ii) \emph{The switched element}, divided by $C\eta'^\theta$, obeys the bound of its type weakened by $\ell_j^{-\theta}$ at its scale (or unweakened if it is fixed): a line with $n\le2$ derivatives by \Cref{rot:h1s:P}, gain $(\eta'/\ell_j)^{1/4}\ge(\eta'/\ell_j)^\theta$ (top piece: $\eta'^{1/4}$); a compound by \Cref{rot:h1s:U}(ii),(iii), gain $\eta'^{\theta_U}\ell_{j_f}^{-\theta_U}$ with $\theta_U\ge\theta$; a first-order compound or a bracket by \Cref{rot:h1s:U}(iv), gain $\ge(\eta'/\ell_j)^{1/16}$ resp.\ $\eta^{1/4}$ (a bracket whose differenced factor is a compound: $\eta'^{\theta_U}\ell_{j_f}^{-\theta_U}$ with $\theta_U\ge\min\{\theta,s/8\}$); a coarse element by $C\eta^{1/2}$; an external by $C\eta^{1/2}\|E\|$, its gradient as a recipient by $C\eta^\kappa\|E\|$ (\Cref{rot:h1s:E}); a vertex function by $c_\tau\eta'\|\nabla w\|$, and the $\partial^{\eta'}g$ term of (R2a), of size $O(\eta')$ in both flows, by $O(\eta')\le\eta'^\theta\ell_{j(v)}^{1-\theta}$; at $e>0$ a contour element by \Cref{rot:prop:BD}(b): $(\eta'/\ell_j)^{1/2}$ for contour lines, also with the one lattice derivative a recipient contour line carries, $\eta^{1/2}$ for contour vertices, $e\,p_r\lmu^{-1}c_\tau\delta_r^{-1/2}\eta'$ for the phases, whence the prefactor; contour vertices are never recipients (\Cref{rot:h1s:rule}(iii)). For $j\le2$ there is no gain and the factor is paid by the degree loss $\theta$, as in \Cref{rot:cor:unif}(a).
(iii) \emph{Degrees.} By \Cref{rot:h1s:Dp} every component of every initial segment of the reduced ordering has $D\ge1$, or $D\in\alpha+\N_0$, or is an explicit-$\eta'$ component with $D=1-\theta$ after $\eta'^\theta$ is taken out; fixed elements impose no condition. Hence $D(H)-\theta>0$ for every component containing the switched element, and the scale sums of \Cref{rot:h1s:MS} converge geometrically there.
(iv) \emph{Logarithms.} $n_0\le\bar n$ zero-degree components, $(1+r)$ per compound, $(1+r)^{n_0+1}$ per common number, $(1+r)^2$ per double compound, and the $(1+r)^{n_0}$ of \Cref{rot:h1s:MS} itself; with $\le\bar n$ elements, $p(\bar n)\le3\bar n+2$ suffices, and $(1+r)\le(1+k)$.
\Cref{rot:h1s:MS} then gives $|T_i|\le C(1+k)^p\prod_lc_l$ with $\prod_lc_l=\eta'^\theta\cdot[\text{per-flow constants}]$; the localisation \cite[(3.55)]{K1} costs at most $p_r^3$ per coarse line, absorbed in the field norms of step (5) of the proof of \Cref{rot:thm:A}.
\end{proof}

\paragraph{The exponent $\theta_1$.}
In every term the switched element carries $\eta'^{\theta'}\ell_{j}^{-\theta'}$ with its own exponent $\theta'$, and \Cref{rot:h1s:MS} at the stage where the component $H\ni e_i$ is shrunk sums $\sum_{j\le j'}\ell_j^{D(H)-\theta}\eta'^\theta\le C\ell_{j'}^{D(H)-\theta}\eta'^\theta$, geometric because $D(H)-\theta>0$; the remaining factor is the merged vertex's weight, and since all later components containing it again have $D-\theta>0$ the factor $\eta'^\theta$ survives. (When $D(H)=0$ the sum $\sum_j(\eta'/\ell_j)^\theta=\sum_j5^{-j\theta}$ is $O(1)$ and the gain is lost; \Cref{rot:h1s:Dp} excludes this.) With $s=\alpha/(1+\alpha)$, $\kappa=\alpha/(2(1+\alpha))$, the gains are: a switched line with $n\le2$ derivatives, $1/2^n\ge\frac14$; the top piece, $\frac14$; a compound with a slice leg carrying $a\le2$ derivatives ($a=2$ only for the differenced compound of (R2b)), $s/2^{a+1}\ge s/8$; with an external or coarse leg, $\kappa s$; the double compound, $s/8$ (slice middle) or $\kappa s$ (coarse middle); the first-order forms, $1/2^{a+2}\ge\frac1{16}$ (slice) or $\frac14$ (external); coarse elements, $\frac12$; vertex functions and explicit-$\eta'$ terms, $1$; contour lines, $\frac12$; contour vertices, $\frac12$; phases, $1$ with the prefactor. The arithmetic: $\kappa s=\alpha^2/(2(1+\alpha)^2)=2\theta_1$; $s/4-\theta_1=\alpha/(4(1+\alpha)^2)>0$; $s/8-\theta_1=\alpha(1-\alpha)/(8(1+\alpha)^2)>0$ for $\alpha<1$; $\frac1{16}-\theta_1=(1-\alpha)(1+3\alpha)/(16(1+\alpha)^2)>0$; and $\min(\alpha,\frac12)-\theta_1>0$ because $\theta_1<\alpha/4\le\alpha$ and $\theta_1<\frac1{16}$. So $\theta_1\defeq\alpha^2/(4(1+\alpha)^2)$ is at most every gain exponent (the weakest two being $s/8$ and $2\theta_1$), and every component containing the switched element keeps the margin $D-\theta_1\ge\min(\alpha,1)-\theta_1>0$. Nothing is nested: no compound is the leg of another (the double compound is one element, and the recipients of (R2) receive a single gradient), so $\theta_1$ does not depend on $\bar n$. The consumers (\Cref{rot:lem:r}, \Cref{rot:prop:cp}) use only $\theta_1>0$ and a polynomial $(1+k)^p$, so the exact value plays no role; what matters is that the set of element types and derivative counts is finite, that every type has a difference law with some positive exponent, and that the local terms are common numbers.

\begin{proof}[Proof of \Cref{rot:thm:H1S}]
Sum \Cref{rot:h1s:T} over the $\le n_\cR$ terms of each ordering class and over the finitely many classes (\cite[(2.7)]{Ba3}; their number depends on the family only). The sectors:
\emph{$e>0$.} The constant-background reduction (\Cref{rot:lem:CB}) is flow-independent algebra; after it the new elements are the contour elements of \Cref{rot:prop:BD}(b). Every connected subgraph containing a contour vertex has degree $\ge\frac12$ (\cite[(3.30)]{K1}; \Cref{rot:lem:CB}(3)), so no new structure of degree $\le0$ appears and \Cref{rot:h1s:Dp} is unchanged; a contour line may be the recipient of (R2a)(ii) at its bond end, and this is the one-derivative clause of \Cref{rot:prop:BD}(b); the phases are Lipschitz in the vertex position with the bound $e\,p_r\lmu^{-1}\delta_r^{-1/2}$, inside the prefactor. $e$ enters only through $e_r^2=e^2\delta_r$, so the bound holds at every $e>0$.
\emph{The sectors with legs at $A=0$.} Scalar two-point: orders $\lambda$, $\lambda^2$, $e^2$, $e^4$ with the counterterm cancellations of \Cref{rot:h1s:VP}(d), $e^2\lambda$ trivial; the $(2,1)$-sector at order $e^3$: T and SC as first-order compounds (no Ward identity of the vertex class is used: in the per-flow normal form the identities of this sector are identities on all-$C$ expressions, applied in each flow alone, and a reduction of this sector to a ``mixed zero-moment'' bound for current--current subgraphs with lines from different flows is not used, the mixed zero moment being a lattice-scale quantity of order $e_r^2/\eta'$); the vacuum polarisation at $e^2$ through \Cref{rot:h1s:VP}(a),(b) and at $e^4$ through (c); L24; everything else of positive degree (\cite[Cor.~2.3, (2.17)]{Ba3}).
\emph{The vacuum sector.} All-$C$ vacuum monomials cancel against $E_1$ exactly (\Cref{rot:lem:ctsame}, \cite[(3.35)]{K1}); the $\Gamma$-monomials, cut at their $\Gamma$-lines, leave all-$C$ components with at least two legs, renormalised by the rule with the $\Gamma$-lines as legs and recipients (one-point components are common numbers, (R2c)); \Cref{rot:h1s:T} applies per coarse site, times $|T|\delta_r^{-3}$ sites (\Cref{rot:thm:V}(e)).
\emph{The source vertex.} The $h$-vertex has degree $2$ (\Cref{rot:lem:S1}); the $h$-families other than $\cT_h$ contain no new structure of degree $\le0$.
Hence $|E^X(\cR)-E^Y(\cR)|\le Cn_\cR\,p_r^2\delta_r^{-1/2}\eta'^\theta(1+k)^p\cdot[\text{right side of \cite[(1.33)]{Ba3}}]$ for $0<\theta\le\theta_1$, uniformly in $r$, $k$, the torus (the tori enter through periodisation, \Cref{rot:h1s:Z0}) and the masses, and in every term of the telescoping the switched element obeys the bound of its type weakened by $\theta$, which is the statement.
\end{proof}
